% s00-front.tex —— 前置：摘要与译本说明

\section*{摘要}
\addcontentsline{toc}{section}{摘要}

受非分歧整体几何 Langlands 对应近期获证的激励，本文朝构造几何 Langlands 对应的驯顺分歧版本推进一步。
我们证明：光滑、射影、连通曲线上，带一点处 Iwahori 层级结构的主 $G$-丛模叠上，
D-模范畴是紧生成的；有了这一点，便可在此情形下构造 Langlands 函子。
随后我们研究该对应的局部行为：在一点处构造因子化模范畴，并在其间构造等价。
这些等价分别是 Arkhipov–Bezrukavnikov 等价、Bezrukavnikov 等价，
以及一个 Iwahori 不变的基本局部等价的因子化模范畴版本。

\section*{译本说明}
\addcontentsline{toc}{section}{译本说明}

\subsection*{一、底本}

\begin{itemize}
\item Taeuk Nam, \emph{Towards a Tamely Ramified Geometric Langlands Correspondence},
哈佛大学文理研究生院博士学位论文，2026 年 4 月，导师 Dennis Gaitsgory，127 页正文。
\item 底本取自哈佛大学 DASH 仓储：\url{https://dash.harvard.edu/handle/1/42744142}。
\item 底本 PDF 共 137 页，其中正文印刷页 $1$–$127$ 对应 PDF 第 $10$–$136$ 页（\textbf{印刷页 $=$ PDF 页 $- 9$}）。
译文正文中的边注 \textsf{p.$N$} 标出该段落在底本中的起始页，便于逐页核对。
\item 底本的扉页、版权页与目录为罗马数字页码，本译本不另立页码，只把内容并入前置件。
\end{itemize}

\subsection*{二、体例}

\begin{itemize}
\item 章、节、小节与定理的编号与底本完全一致。底本的定理编号形如 \textbf{2.3.1.2}
（章.节.小节.序号），译本沿用同一四级结构，故编号可直接对照；公式编号形如 \textbf{(2.3.1)}，
亦与底本一致。
\item 译者注以脚注形式排于页脚，楷体，并标注来源等级：
A 为原文、作者本人论文、经典教材或 arXiv 预印本；B 为 nLab、Stacks Project、ProofWiki 等白名单站点
（须与 A 类交叉核实）；I 为本译丛既有译注、术语表与资料卡。无外部依据者标「译者按（无外部依据）」。
\item 底本明显的笔误，译本一律照排原文并加注说明，不径改正文；全部条目登记于
\texttt{工程文件/work/疑点清单.tsv}。
\item 术语依 \texttt{工程文件/work/glossary.tsv} 冻结；未定名者保留原文并加注。
\end{itemize}

\subsection*{三、术语约定}

本译本沿用本译丛在几何 Langlands 方向已冻结的译名，主要有：
Iwahori 子群、Whittaker 范畴、球面 Hecke 范畴、Iwahori–Hecke 范畴、oper（保留原词）、
Ran 空间、因子化范畴、可熔（fusable）、临界层级、紧生成、截断的 / 可截断的 / 收缩的。
$G$ 的对偶群记 $\Gd$，$\check{\mathfrak{n}}/\check B$ 记 $\nd/\Bd$。
完整对照见术语表；同一术语在全书内不得改译。

\subsection*{四、分工与责任}

本译本由数译工程组织完成：底本逐页切分后分章精读与翻译，再由统一体例合成、逐页核对公式，
并逐条登记疑点。译文对底本的忠实度以「不改原文、只加批注」为限；
读者若发现译文与底本图像不符，请以底本为准并告知我们修订。
